\section{Problem Formulation}
\label{sec:problem}

% PF-P1 — Workload definition
We consider a complex information need $q$ over one or more evidence sources
$\mathcal{D}$ under a finite resource budget $B$. Solving $q$ requires
acquiring a set of mutually dependent evidence items before producing an answer
(or, for verification tasks, a verdict). The system therefore optimizes not
only answer generation but the construction of an \emph{evidence state} that
satisfies the information requirements induced by $q$. Two properties make the
problem interesting rather than a single lookup. First, the evidence items are
\emph{dependent}: binding one item (e.g., an entity) is a precondition for
acquiring another (e.g., the relation that connects it to a second entity).
Second, the system operates under \emph{partial information}: whether a bridge
binding exists, how many candidates a retrieval step will return, and whether
the answer model can use the materialized evidence are all uncertain before
execution. A plan must therefore commit resources before the outcomes that
determine its success are observed.

% PF-P2 — Evidence item
An evidence item is represented as a typed tuple
$e = \langle \mathit{id}, \mathit{type}, \mathit{value}, \mathit{source},
\mathit{provenance}, \mathit{bindings}, \mathit{cost}, \mathit{quality}
\rangle$ (\code{EvidenceRecord}). Unlike an untyped text chunk, this
representation exposes the properties downstream operators and the optimizer
need: the \emph{type} determines which operators can consume it, the
\emph{source} and \emph{provenance} make it traceable, the \emph{bindings}
record which variables it instantiates, and \emph{cost}/\emph{quality} let the
optimizer reason about resource consumption and expected utility. Treating an
evidence item as a first-class record rather than a passage string is what
makes typed composition and requirement-aware optimization well-defined.

% PF-P3 — Evidence requirement
An evidence requirement $r$ specifies a condition that must be satisfied rather
than an action that must be executed (\code{EvidenceRequirement}): a set of
variables to bind, an expected evidence type, an importance, and a status
(\code{unresolved}/\code{partial}/\code{satisfied}). Requirements may depend on
bindings produced by other requirements; this dependency is expressed as a
relation over requirement ids and induces a dependency graph that
\emph{constrains}---but does not fully determine---the legal execution order.
The distinction between \emph{condition} and \emph{action} is the crux of the
paper's positioning: the information need is stable, while the physical actions
used to satisfy it are revisable and budget-bound.

% PF-P4 — Evidence state
Execution evolves an evidence state $S_t$ (\code{EvidenceState}) containing the
acquired evidence, the current bindings, each requirement's status, the
execution history, and the realized resource usage. Because retrieval success,
selectivity, and evidence quality are only partially observable before
execution, each state transition can in principle change both the feasible
action space and the preferred physical plan. In the current system the
execution topology is chain-only, so the feasible action space is a singleton
order and state transitions change bindings but not the order; the formal
object is nonetheless the same one over which the branching-execution extension
would operate (Section~\ref{sec:execution}).

% PF-P5 — Optimization objective
We formulate physical planning as constrained optimization: maximize expected
satisfaction of weighted evidence requirements and answer-facing evidence
utility subject to explicit retrieval, context-token, latency, and optional
monetary budgets. The objective (Equation~\ref{eq:objective}) separates task
utility from operational constraints. Utility is credited only for progress on
unresolved requirements; a requirement's importance weight is derived from its
dependency depth by the chain law ($\tau = 2d-1$) on well-defined chains. This
formulation enables \emph{matched-budget} comparison: two plans are compared at
the same budget ceiling, and the realized cost dimension is reported alongside
quality rather than hidden.

% PF-P6 — Why optimization is non-trivial
The optimization problem is combinatorial because each logical operator may
admit multiple physical implementations, dependent operators have partially
constrained orderings, and budget must be allocated before uncertain evidence
yields are observed. Concretely, the search space is the product of legal
orders, per-requirement physical implementations, and budget allocations.
Section~\ref{sec:optimizer} derives the resulting search problem and motivates
the allocation procedure: on the well-defined chains that dominate our
workloads the dependency graph admits exactly one legal order, so the
procedure reduces to an importance-weighted ($\tau = 2d-1$) budget allocation
over a fixed order; on general graphs it falls back to bounded enumeration with
dominance pruning. The non-triviality is not merely
computational: because yields are uncertain, a plan that spends more on a
high-importance downstream requirement can be strictly better than a cheaper
plan even when the cheaper plan is closer to the budget ceiling---so a
cost-only or relevance-only objective is provably insufficient.
